\subsection{Elimination of the entire three-primary reduced obstruction}

\begin{result}{B.2.23}{Proved}{The remaining central obstruction is two-primary}
For the actual universal central kernel and its two
torsion correction maps one has
\[
 M_{(3)}=\operatorname{im}\kappa_r\cong C_3,
                    \qquad (M/S)_{(3)}=0.                    \tag{B.S55}
\]
Consequently $\overline M=M/S$ is a finite two-group
generated by at most four elements and annihilated
by $2^{14}$.  In particular
\[
 |\overline M|\leq 2^{56},
                   \qquad [H:D_P(H)]\leq 2^{56}.              \tag{B.S56}
\]
The second inequality is a cardinality bound for
the affine orbit in B.2.19; no assertion that its
cardinality is a power of two is needed.
\end{result}

\begin{proof}
Use the specified ternary detector $R_3$ of B.2.22.
Its normal two-core has order $2048=2^{11}$.
Its map to $E$ identifies this core with $O_2(E)$:
the latter also has order $2048$, and the kernel
$A_3^6$ is a three-group.  The actual quotient maps
therefore give a commutative diagram of extensions
\[
\begin{array}{ccccccccc}
 1&\longrightarrow&A_3^6&\longrightarrow&R_3&\longrightarrow&E&\longrightarrow&1\\
 &&\Vert&&\downarrow&&\downarrow\\
 1&\longrightarrow&A_3^6&\longrightarrow&\widetilde R_3&\longrightarrow&\widetilde E&\longrightarrow&1,
\end{array}                                                     \tag{B.S57}
\]
where
\[
 |\widetilde R_3|=2^3 3^{30},
       \qquad |\widetilde E|=2^3 3^6=5832,
       \qquad A_3^6=C_3^{24}.                                 \tag{B.S58}
\]
All four groups and maps in (B.S57) are marked by
the same images of $s,r$.  Quotienting a normal
group of order prime to three induces an
isomorphism on cohomology with trivial
$\mathbb F_3$ coefficients, by averaging and the
quotient spectral sequence.  Thus this reduction
preserves both the cohomology of $R_3$ and its
restriction to the displayed $A_3^6$.

Here the complete finite calculation is
\[
\begin{array}{c|c}
 \text{quantity}&\text{value}\\ \hline
 \dim H^2(\widetilde R_3,\mathbb F_3)&3\\
 \text{rank of restriction to the source torsion factors}&1\\
 \text{dimension of the marked central layer}&2\\
 \text{rank of restriction to }A_3^6&1
\end{array}                                                     \tag{B.S59}
\]
The source restriction vector in the computed
cohomology basis is $(1,2,1)$.  After taking its
kernel, the marked central cover has kernel
$C_3^2$.  It is generated by the named torsion
lifts, and the $r$-centralizer correction has
image of order three.  These calculations are
performed in
\path{work/centralizer_code/check_universal_small_homology_3_reduced.g}
and its executed output; the exact power--conjugate
presentation and both cocycles are exported in
\path{work/centralizer_code/small_homology_cover_3_reduced.json}.

The final restriction rank is evaluated independently
from those exported group and cocycle data in
\path{work/centralizer_code/check_ternary_restriction.g}.
The script constructs the actual homomorphism to
$\widetilde E$ by $s\mapsto s$, $r\mapsto r$, and
checks that its kernel is elementary abelian of
order $3^{24}$.  It then evaluates the cubes and
all pairwise commutators of lifts of a basis of
this kernel.  In the two central coordinates the
only nonzero rows are $(0,1)$ and $(0,2)$, so their
span has rank one.  For an elementary abelian
three-group these cubes and commutators determine
a central extension class completely: zero rows
give an abelian extension of exponent three,
which splits as a vector-space sequence.
Thus these word evaluations give precisely the
last rank in (B.S59).

The missing one-dimensional source class also
restricts to zero on $A_3^6$.  It is the
universal-coefficient class from the abelianization
$C_6$, and $A_3^6$ lies in the derived subgroup,
since $\widetilde R_3^{\mathrm{ab}}\to
\widetilde E^{\mathrm{ab}}$ is the isomorphism
$C_6\to C_6$.  The rank one in (B.S59) is
therefore the rank of restriction on the entire
three-dimensional second cohomology, not only
on the marked subspace.

In the cohomology spectral sequence of
$A_3^6\to R_3\to E$, the first filtration step
of total degree two is the kernel of restriction
to $A_3^6$.  It has dimension $3-1=2$ by (B.S59).
By (B.S51)--(B.S52), the comparison with
$L^6\to U_0\to E$ is an isomorphism on this
filtration step: its two graded pieces in
bidegrees $(2,0)$ and $(1,1)$ have identical
incoming and outgoing differentials, as checked
in the proof of (B.S53).  The infinite group's
degree-two invariant term from $L^6$ is zero
at three by (B.S42).  Its entire second
cohomology is consequently this first filtration
step.  We obtain the exact dimension
\[
                    \dim H^2(U_0,\mathbb F_3)=2.              \tag{B.S60}
\]

The universal-coefficient sequence subtracts one
dimension for $H_1(U_0,\mathbb Z)=C_6$, leaving
one generator for the three-primary part of
$H_2(U_0,\mathbb Z)$.  Its quotient $M_{(3)}$
therefore is cyclic.  The splitting in B.2.20
already gives an order-three subgroup
$\operatorname{im}\kappa_r\subseteq M_{(3)}$
which maps onto $(M_E)_{(3)}=C_3$.  In a cyclic
group of order $3^a$ with $a>1$, its subgroup of
order three maps trivially under every quotient
of order three.  Here that map is onto.
Consequently $a=1$, proving (B.S55).

The remaining assertions now follow from B.2.21
and the exact affine-orbit formula of B.2.19.
\end{proof}
